% Created 2026-09-30 Wed 15:28
% Intended LaTeX compiler: lualatex
\documentclass[11pt]{article}

\usepackage{amsmath}
\usepackage{fontspec}
\usepackage{graphicx}
\usepackage{longtable}
\usepackage{wrapfig}
\usepackage{rotating}
\usepackage[normalem]{ulem}
\usepackage{capt-of}
\usepackage{hyperref}
\usepackage{fontspec}
\setmonofont{DejaVu Sans Mono}[Scale=MatchLowercase]
\date{\today}
\title{Debugging Python with dape}
\hypersetup{
 pdfauthor={},
 pdftitle={Debugging Python with dape},
 pdfkeywords={},
 pdfsubject={},
 pdfcreator={},
 pdflang={English}}
\usepackage{biblatex}

\begin{document}

\maketitle
How the dape setup in \texttt{config.el} looks and works.  The keys follow PyCharm,
except where KDE takes Ctrl+F1..F12 for switching desktops.
\section{Keys}
\label{sec:org64fe645}

\begin{center}
\begin{tabular}{lll}
Action & PyCharm & Here\\
\hline
Start debugging & \texttt{Shift+F9} & \texttt{S-<f9>} (then \texttt{debugpy})\\
Resume & \texttt{F9} & \texttt{<f9>}\\
Step Over & \texttt{F8}     >\textasciitilde{}\\
Step Into & \texttt{F7} & \texttt{<f7>}\\
Step Out & \texttt{Shift+F8} & \texttt{S-<f8>}\\
Toggle breakpoint & \texttt{Ctrl+F8} / click gutter & \texttt{C-S-b} / left-click fringe\\
Stop & \texttt{Ctrl+F2} & \texttt{S-<f5>}\\
Conditional breakpoint & right-click breakpoint & middle-click fringe / \texttt{C-x C-a e}\\
Log breakpoint (no suspend) & \texttt{Shift}-click gutter & right-click fringe / \texttt{C-x C-a l}\\
Run to Cursor & \texttt{Alt+F9} & \texttt{C-x C-a u}\\
Evaluate Expression & \texttt{Alt+F8} & \texttt{C-x C-a x}, or type it in the REPL\\
Add to Watches & right-click → Add to Watches & \texttt{C-x C-a w} (symbol at point)\\
Rerun & \texttt{Ctrl+F5} & \texttt{C-x C-a r}\\
Pause & Pause button & \texttt{C-x C-a p}\\
Info panels (close/reopen) & Debug tool window & \texttt{C-x C-a i}\\
Debug console & Console tab & \texttt{C-x C-a R} (REPL)\\
Up/down the call stack & click in Frames & \texttt{C-x C-a <} / \texttt{C-x C-a >}\\
Remove all breakpoints & Remove All Breakpoints & \texttt{C-x C-a B}\\
\end{tabular}
\end{center}

\texttt{C-<f8>} and \texttt{C-<f2>} are bound too, but only reach Emacs if KDE's ``Switch to
Desktop 8/2'' shortcuts are freed.  The \texttt{C-x C-a} prefix exists once dape has
loaded, which happens when a Python file is opened.  Not in dape: Smart Step
Into, Mute Breakpoints, Force Step Into.
\section{The screen while debugging}
\label{sec:orga28fd0d}

\begin{verbatim}
┌──────────┬─────────────────────────────────┬────────┐
│ Treemacs │ ●  41  def step(x):             │ Claude │
│          │ ▶  42      y = x * 2    x: 3    │        │
│──────────│    43      return y + 1         │        │
│ Outline  ├──────────┬──────────┬───────────┤        │
│          │ Stack    │ Scope    │ REPL      │        │
│          │ (TAB ↻)  │ (TAB ↻)  │           │        │
└──────────┴──────────┴──────────┴───────────┴────────┘
\end{verbatim}

\begin{itemize}
\item In the editor, \texttt{●} in the fringe (the thin strip at the window's left
edge, left of the line numbers) is a breakpoint, \texttt{▶} marks the line about
to run, and the grey text after the current and previous lines shows the
values of their variables.
\item The terminal steps aside while a session runs and comes back on quit.
\item Each panel's header line lists its tabs; \texttt{TAB} / \texttt{S-TAB} cycle through them,
or click one.  Left panel: Stack, Breakpoints, Threads, Modules, Sources.
Middle panel: Scope, Watch.
\end{itemize}
\section{A first session}
\label{sec:orgb7058dc}

\begin{enumerate}
\item Open a script and left-click the fringe next to a line inside a function
or loop (or \texttt{C-S-b} on that line).
\item \texttt{S-<f9>}.  The first time, type \texttt{debugpy} (\texttt{TAB} completes) and \texttt{RET}; it
runs the current file.  Afterwards the prompt comes pre-filled with the
last command.  To pass command-line arguments, add them before \texttt{RET}:
\texttt{debugpy :args ["-{}-{}steps" "100"]}.  (\texttt{debugpy-module} runs the current
directory as a package, like \texttt{python -m}.)
\item The program runs up to the breakpoint; Scope shows the local variables.
\item \texttt{<f8>} a few times to step over lines and watch the values change.
\texttt{<f7>} on a function call steps into it; \texttt{S-<f8>} steps back out.
\item In Stack, move to a lower frame and \texttt{RET}: the editor jumps to the caller
and Scope shows the caller's variables.  Stepping always continues from
the top frame.
\item \texttt{<f9>} continues to the next breakpoint.
\item \texttt{S-<f5>} quits: the panels close and the terminal comes back.  If the
program finishes by itself, the REPL says ``Session terminated''; \texttt{S-<f5>}
still closes the panels.
\end{enumerate}
\section{The panels}
\label{sec:orgdf461ba}

\begin{center}
\begin{tabular}{lll}
Panel & Key & What it does\\
\hline
Stack & \texttt{RET} & Select a frame: jump to its code and show its variables\\
Breakpoints & \texttt{RET} & Go to the breakpoint; on the exception rows (raised / uncaught)\\
 &  & turn breaking on exceptions on or off\\
 & \texttt{D} / \texttt{d} & Disable a breakpoint / delete it\\
 & \texttt{e} & Edit a log breakpoint's message\\
Scope & \texttt{e} & Expand or collapse a dict, list or object\\
 & \texttt{W} & Add the variable to Watch\\
 & \texttt{=} & Change the variable's value while paused\\
Watch & \texttt{W} & Remove the expression from Watch\\
 & \texttt{C-x C-q} & Edit the list of watch expressions as text\\
\end{tabular}
\end{center}

In Scope and Watch, \texttt{RET} does one of those three things depending on where
the cursor is: on the \texttt{+} / \texttt{-} it expands, on the name it adds to Watch, on
the value it edits.  To watch a variable straight from the code, put the
cursor on it and press \texttt{C-x C-a w}.
\section{The REPL}
\label{sec:orga355352}

\begin{itemize}
\item It shows everything the program prints.  The program can't read keyboard
input (\texttt{input()}) in this setup.
\item While paused, type any Python expression and it is evaluated in the
selected frame: \texttt{len(sensors)}, \texttt{x * 2}, \texttt{vars(self)}.  Assignments such as
\texttt{x = 5} work too.
\item One-letter commands: \texttt{n} step over, \texttt{s} step into, \texttt{o} step out,
\texttt{c} continue, \texttt{r} restart, \texttt{q} quit.  \texttt{RET} on an empty line repeats the
last command, so \texttt{n} \texttt{RET RET RET} steps quickly.
\item A variable named like a command (\texttt{b c d e k m n o p q r s t u w}) runs the
command instead.  Type a space before the name to see the variable, as in `` n''.
\end{itemize}
\section{Breakpoints that don't just stop}
\label{sec:orgf75866a}

\begin{itemize}
\item Conditional: middle-click the fringe (or \texttt{C-x C-a e}) and enter a
condition such as \texttt{i == 50}.  The program stops only when it is true.
\item Log: right-click the fringe (or \texttt{C-x C-a l}) and enter a message such as
\texttt{step \{i\}: x=\{x\}}.  The program doesn't stop; each hit prints the message
in the REPL.  Expressions go in \texttt{\{\}}.
\end{itemize}
\end{document}
